\documentclass[11pt,a4paper]{article}

\usepackage[margin=2.4cm]{geometry}
\usepackage{enumitem}
\usepackage{parskip}
\usepackage{titlesec}

\titleformat{\section}{\large\bfseries}{\thesection.}{0.6em}{}
\titleformat{\subsection}{\normalsize\bfseries}{\thesubsection}{0.6em}{}
\setlist[description]{style=nextline, leftmargin=1.2em, itemsep=0.4em}
\newcommand{\code}[1]{\texttt{#1}}

\title{\vspace{-1.5cm}Crypto Core (RUB) --- Boot-Stage Scoping Brief}
\author{Benjamin Schulte}
\date{\today}

\begin{document}
\maketitle
\vspace{-0.6cm}

\noindent\textit{Purpose: align on the scope and ownership decisions needed to implement the boot
stage (secure key provisioning), after which the crypto core is implementable up to chip
integration. The runtime datapath is complete and verified; the open items are security-scope
and system-integration choices that the core cannot decide in isolation.}

%==================================================================
\section{What the core can do today}
%==================================================================

\begin{description}
  \item[Interfaces] Two AXI ports, matching the SoC accelerator pattern:
    a \textbf{REG-IF} (AXI4-Lite \emph{slave}) for the host to configure the core, and a
    \textbf{DATA-IF} (AXI4 \emph{master}) over which the core autonomously moves data
    to/from the crossbar (a ``crypto-DMA'').
  \item[Cipher \& modes] AES-128 in \textbf{CTR mode} (encryption and decryption are the same
    operation). Two independent key slots with runtime selection (\code{key\_sel}), and a
    \textbf{bypass} mode for plain data movement (no crypto).
  \item[Autonomous operation] The host programs source/destination addresses and a block count,
    sets \code{GO}, and the core streams: read burst $\rightarrow$ AES-CTR $\rightarrow$ write burst,
    without per-block host involvement.
  \item[Nonce-uniqueness guarantee] 96-bit host-supplied IV + 32-bit hardware counter that
    auto-increments per block and cannot be rolled back. A pre-flight gate \textbf{refuses} any
    transfer that would exceed $2^{32}$ blocks on one IV (which would reuse a nonce) and raises a
    \code{dma\_error} status bit; the host recovers by writing a fresh IV.
  \item[Hardening] Keys are never readable over the bus (no AXI read path to key storage);
    keys are asynchronously zeroized on reset.
  \item[Verification] Self-checking testbenches all pass: functional (NIST SP800-38A CTR vectors,
    multi-block, encrypt/decrypt round-trip, bypass copy) and safety (counter-overflow boundary,
    key isolation, reset-during-operation, no-key refusal), plus unit tests for the register file,
    FSM, and key store.
\end{description}

\noindent\textbf{Not yet implemented:} secure key provisioning (keys are currently loaded only via a
simulation-only mechanism), and any form of \emph{authentication / integrity} --- the core today is
confidentiality-only.

%==================================================================
\section{Scope decisions (security level)}
%==================================================================

\noindent These decide how large the security story is. They should be answered before the boot
stage is detailed, because they change what hardware we must build.

\begin{description}

  \item[(0) Threat model \& compliance target --- the framing decision.]
    What are we defending against: untrusted external flash, untrusted host software, a physical
    attacker (side channels)? Is there a compliance target (FIPS, Common Criteria, internal
    standard)? Every decision below scales with the answer.

  \item[(1) Do we want authenticated \emph{encryption} on the data path?]
    Today CTR provides confidentiality only --- tampered ciphertext is undetectable.
    \emph{If yes} $\rightarrow$ we add \textbf{AES-GCM}, which requires a GF($2^{128}$) multiplier
    (reuses the existing AES engine).
    \emph{If no} $\rightarrow$ we stay CTR-only and accept that data integrity is a host/protocol
    responsibility.

  \item[(2) Do we want an authenticated / signed \emph{bitstream} (firmware integrity)?]
    \emph{If yes} $\rightarrow$ choose the mechanism:
    \textbf{symmetric} (CMAC/GMAC with a device key --- cheap if GCM already exists), or
    \textbf{asymmetric} (a digital signature, e.g.\ ECDSA/RSA --- needs a public-key verifier and a
    trust anchor, a substantial separate block). Also decide \emph{whose} block this is: the crypto
    core, or a dedicated secure-boot unit in the SoC.
    \emph{If no} $\rightarrow$ a tampered bitstream is not detected.

  \item[(3) Do we want authenticated \emph{key wrapping} at boot (tamper-proof keys)?]
    A plain CTR-unwrap of the keys is malleable: an attacker who can edit the stored blob can flip
    key bits deterministically.
    \emph{If yes} $\rightarrow$ authenticated unwrap (AES-GCM or NIST AES Key-Wrap KW/KWP).
    \emph{If no} $\rightarrow$ confidentiality-only key unwrap, vulnerable to targeted bit-flips.

  \item[(4) Is side-channel resistance in scope?]
    \emph{If yes} $\rightarrow$ AES masking (our AES core may already expose a security/area
    parameter worth checking) and DPA-resistant key handling.
    \emph{If no} $\rightarrow$ no countermeasures, suitable only for a non-physical threat model.

\end{description}

\noindent\textbf{Key dependency:} decisions (1), (2), and (3) all share \textbf{one} building block ---
an authenticated mode (AES-GCM / GMAC). If any of them is ``yes'', it is most efficient to build GCM
once and reuse it for the data path, the bitstream MAC, and the key unwrap. This is the single most
consequential outcome to settle in the meeting.

%==================================================================
\section{Boot stage: root key, format \& ownership (who does what)}
%==================================================================

\noindent Assuming the boot stage proceeds, the following must be specified --- several are
shared between us and the SoC / manufacturing teams.

\begin{description}

  \item[(A) Root key (KEK) --- source and delivery.]
    Where does the device root key come from: \textbf{eFuse}, \textbf{PUF}, or a tie-off for the
    test chip? Is it \emph{per-device unique} or shared across a batch (a shared KEK with a fixed
    boot nonce reuses the same unwrap keystream on every device)? Who provisions it, and at which
    lifecycle stage? It must reach the core as an \emph{internal hardware input} and never be
    exposed on the AXI bus.

  \item[(B) Wrapped-key blob --- location and retrieval.]
    Where is the encrypted key material stored: embedded in the bitstream, in external
    \textbf{QSPI flash}, or on-chip ROM? How does the core obtain it: \emph{fetch it itself over the
    DATA-IF} (reusing the AXI master), receive it on \emph{dedicated ports}, or have a \emph{system
    boot controller} hand it over? Is boot \emph{autonomous at power-on} or orchestrated by an SoC
    boot sequencer that releases the core?

  \item[(C) Blob / bitstream layout --- format and ownership.]
    Who owns and writes the format specification? What fields does it carry: header/magic, version,
    boot nonce/IV, wrapped key\,1, wrapped key\,2, and (if decision~3 is ``yes'') an authentication
    tag? What word ordering / endianness (to match the conventions already used on the register and
    AXI interfaces)? Is the key count fixed at two, or variable?

  \item[(D) Boot failure \& lifecycle.]
    On failure (wrong KEK or failed authentication): zeroize, lock the core, set an error status,
    halt --- confirm the desired behavior. Is field \emph{re-provisioning / key rotation} required?
    Is there a \emph{production lifecycle state} that must disable the lab key-load path (secure vs.\
    open device states)?

\end{description}

%==================================================================
\section{Integration parameters to finalize}
%==================================================================

\begin{description}
  \item[Bus widths] AXI Crossbar data and address width --- these fix the core's
    \code{AXI\_DATA\_WIDTH} / \code{AXI\_ADDR\_WIDTH} parameters (currently placeholders).
  \item[Notification] An interrupt line for done/error, or polling only (polling works today).
  \item[Mapping \& timing] The REG-IF base address in the system memory map; clock/reset domains
    and target frequency.
\end{description}

\vspace{0.4cm}
\noindent\textit{Recommended priority: settle the framing threat model (0) and the authentication
question (1)--(3) first --- they determine whether GCM is built now --- then the root-key source (A)
and blob retrieval (B), which unblock the boot-FSM design.}

\end{document}
